% Yash Vadi - Resume
\documentclass[letterpaper,10pt]{article}

\usepackage[letterpaper,top=0.38in,bottom=0.38in,left=0.55in,right=0.55in]{geometry}
\usepackage[T1]{fontenc}
\usepackage{lmodern}
\usepackage{microtype}
\usepackage{titlesec}
\usepackage{enumitem}
\usepackage{tabularx}
\usepackage[dvipsnames]{xcolor}
\usepackage[colorlinks=true, urlcolor=NavyBlue, linkcolor=NavyBlue, citecolor=NavyBlue]{hyperref}

% XeTeX/Tectonic already emits Unicode text mappings; keep the source
% compatible with pdfLaTeX when that engine is used instead.
\ifdefined\pdfglyphtounicode
  \input{glyphtounicode}
  \pdfgentounicode=1
\fi

\pagestyle{empty}
\setlength{\parindent}{0pt}
\setlength{\parskip}{0pt}
\urlstyle{same}
\raggedright

\titleformat{\section}{\vspace{3pt}\large\bfseries}{}{0pt}{}[\vspace{1pt}\titlerule]
\titlespacing*{\section}{0pt}{3pt}{2pt}
\setlist[itemize]{leftmargin=13pt,itemsep=0pt,topsep=1pt,parsep=0pt,partopsep=0pt}

\newcommand{\entry}[4]{%
  \textbf{#1}\hfill #2\\[-2pt]
  \textit{#3}\hfill \textit{#4}\vspace{-3pt}
}
\newcommand{\subentry}[2]{\textit{#1}\hfill \textit{#2}\\[-2pt]}
\newcommand{\pub}[1]{\item\small #1}

\begin{document}
\small

%==================== HEADER ====================
\begin{center}
  {\LARGE\textbf{Yash Vadi}}\\[3pt]
  Toronto, ON, Canada \;|\;
  \href{mailto:yash.vadi@yahoo.com}{\textcolor{NavyBlue}{yash.vadi@yahoo.com}} \;|\;
  +1 (438) 924-4257 \;|\;
  \href{https://yashvadi.github.io/}{\textcolor{NavyBlue}{yashvadi.github.io}}
\end{center}
\vspace{-8pt}

%==================== EDUCATION ====================
\section{Education}
\textbf{M.Sc. in Computer Science, Machine Learning}, Universit\'e de Montr\'eal -- MILA \hfill 2023 -- 2025\\
\textbf{B.Tech. in Electronics and Communication Engineering}, Indian Institute of Information Technology, Surat \hfill 2019 -- 2023

%==================== PUBLICATIONS ====================
\section{Selected Publications}
\begin{itemize}
  \pub{S. Joseph et al. (including \textbf{Y. Vadi}). ``Prisma: An Open Source Toolkit for Mechanistic Interpretability in Vision and Video.'' \textit{Mechanistic Interpretability for Vision Workshop, CVPR 2025} (Oral). \href{https://arxiv.org/abs/2504.19475}{[arXiv]}}
  \pub{S. D. Das*, \textbf{Y. Vadi*}, and K. Yadav. ``Transformer-based Joint Modeling for Automatic Essay Scoring and Off-Topic Detection.'' \textit{LREC-COLING 2024}. (* Equal contribution.) \href{https://aclanthology.org/2024.lrec-main.1457/}{[Paper]}}
  \pub{S. D. Das, \textbf{Y. Vadi}, A. Unnam, and K. Yadav. ``Unsupervised Out-of-Distribution Dialect Detection with Mahalanobis Distance.'' \textit{INTERSPEECH 2023}. \href{https://www.isca-archive.org/interspeech_2023/das23_interspeech.html}{[Paper]}}
\end{itemize}

%==================== EXPERIENCE ====================
\section{Experience}

\entry{ServiceNow}{Toronto, ON, Canada}{Research Engineer, AI Search}{September 2025 -- Present}
\begin{itemize}
  \item Built a search data-curation and relevance-evaluation framework for enterprise-specific AI Search, using LLM-assisted labeling and A/B benchmarks to assess search quality across datasets.
  \item Evaluated embedding models, cross-encoder rerankers, and alternative retrieval/ranking configurations on search benchmarks to identify effective components for integration into the search stack.
  \item Trained and evaluated L2-stage rerankers to improve candidate quality before results are passed to downstream post-rankers.
\end{itemize}

\entry{ServiceNow}{Montreal, QC, Canada}{Applied Research Scientist}{March 2025 -- September 2025}
\begin{itemize}
  \item Applied supervised fine-tuning (SFT) with LoRA to Llama 3.2 3B, a small language model (SLM), for invoice information extraction; evaluated key-field and line-item extraction performance.
  \item Developed and benchmarked a chunked extraction workflow for long tables in multi-page enterprise documents, using page-level row summaries and bounded chunks to handle context limits; analyzed failures to guide improvements.
  \item Co-authored a patent application for an algorithm to locate multi-line information in complex document layouts.
\end{itemize}

\entry{ServiceNow}{Montreal, QC, Canada}{Applied Research Scientist Intern}{May 2024 -- October 2024}
\begin{itemize}
  \item Developed a synthetic data-generation method for multimodal, layout-aware invoices to increase document diversity for model fine-tuning.
  \item Fine-tuned LILT, a layout-aware multimodal transformer, for PII redaction on invoices to support privacy workflows.
\end{itemize}

\textbf{SHL}\hfill Gurgaon, India\\[-1pt]
\subentry{Research Intern}{January 2023 -- July 2023}
\begin{itemize}
  \item Evaluated AI-generated text detection using GPT-2 perplexity and GLTR token-rank features; published findings on the \href{https://www.shl.com/resources/by-type/blog/2023/the-battle-of-the-pens-insights-into-detecting-ai-generated-text/}{[SHL blog]}.
  \item Fine-tuned encoder-only and encoder--decoder models for personalized follow-up question generation in an AI-based virtual interviewer, and evaluated generated questions using annotated examples.
\end{itemize}
\subentry{Research Intern}{May 2022 -- July 2022}
\begin{itemize}
  \item Implemented and evaluated transformer-based off-topic essay detection; developed an adversarial attack pipeline to assess the robustness of automated essay-scoring models.
\end{itemize}
\subentry{AI Intern}{December 2021 -- March 2022}
\begin{itemize}
  \item Developed a BERT-based intent classifier with Mahalanobis-distance out-of-domain detection (97\% weighted F1 for in-domain classification; 98\% AUROC for out-of-domain detection).
\end{itemize}

%==================== OPEN SOURCE ====================
\section{Open-Source Contributions}
\textbf{ViT-Prisma} \hfill \href{https://github.com/soniajoseph/ViT-Prisma}{[GitHub]}\\[-1pt]
\textit{Open-source toolkit for mechanistic interpretability in Vision Transformers}
\begin{itemize}
  \item Implemented data-parallel, multi-GPU training pipelines for Vision Transformers on classification and patch-reconstruction tasks.
  \item Evaluated model behavior across training configurations and contributed trained checkpoints to the open-source research community.
\end{itemize}

%==================== SKILLS ====================
\section{Technical Skills}
\textbf{Programming:} Python, C, C++\\
\textbf{ML Frameworks:} PyTorch, TensorFlow\\
\textbf{Research Areas:} Machine learning, deep learning, NLP, computer vision, multimodal learning, information retrieval\\
\textbf{Methods:} Model fine-tuning, data curation, experimental design, model evaluation, distributed training

%==================== AWARDS ====================
\section{Awards}
\begin{itemize}
  \item UTG Llama Award, ServiceNow (Q3 2025).
  \item Kaggle Silver Medal, RSNA-MICCAI Brain Tumor Radiogenomic Classification (52nd among 1,500+ participants).
  \item Mind to Market Innovation Challenge: awarded INR 147,499 (approximately USD 2,000) for a drone-based surveillance proposal.
\end{itemize}

\end{document}
